\documentclass[10pt,a4paper,oneside]{article}

\usepackage[utf8]{inputenc}
\usepackage{standalone}
\usepackage{amsmath}
\usepackage{amsfonts}
\usepackage{amssymb}
\usepackage{graphicx}
\usepackage{tocloft}
\usepackage{hyperref}
\usepackage{float}
\usepackage{darkmode}

\renewcommand{\contentsname}{Summary}
\renewcommand{\cftdotsep}{2}
\renewcommand{\cftsecleader}{\cftdotfill{\cftdotsep}}

\enabledarkmode
\setlength{\parindent}{0pt}
\vspace{2cm}

\hypersetup
{
  colorlinks=true,
  linkcolor=black,
  urlcolor=cyan,
  linkcolor=white
}

\author
{
  Pierre Élise \\ \texttt{ISIMA 27, ZZ3 F1}
  \and
  Vérin Clément \\ \texttt{ISIMA 27, ZZ3 F1}
}

\title{ZZ3 : Projet d'Objets Connectés}

\begin{document}

\fbox
{
  \parbox{\textwidth}
  {
    \centering
    \maketitle
  }
}

% \begin{figure}[H]
%   \centering
%   \includegraphics[width=1.0\textwidth]{./LateX/Captures/Title.png}
%   \caption{Selfie of the \textit{Curiosity} Rover (October 2021)}
% \end{figure}

\begin{center}
  \href{https://github.com/ForoyV8-MF/ZZ3-Objets-Connectes-Projet-Pierre-Verin.git} {GitHub Deposit Link}
  \href{https://docs.google.com/document/d/122-0kRPIf3k31qxZWXeNSvJEtt5LDkis1RGscYJJpbM/edit?tab=t.0} {Fichier Partagés Drive}
  \href{https://fr.overleaf.com/project/6abb691e5b82a3c004bbb391} {Projet Overleaf (Outil pour Latex)}
\end{center}

\newpage

\tableofcontents

\newpage

\section{Introduction}

L'objectif est de faire fonctionner un capteur de température qui envoit des données sécurisées avec un modèle personnalisé, et non celui de LoraWan directement.\\

\section{Installation de l'Environnement}

Nous utilisons l'environnement PlatformIO, qui requiert un fichier platformio.ini, 

\section{Bavardage avec la Carte}

La console de Debug représente le terminal intégré de la carte.\\
PlatformIO compile et génère le fichier binaire seulement : il faut bien penser à Uploader.\\

\section{Assignation des Pattes}

Nous cherchons dans l'arborescence un fichier type Board.txt pour l'assignation de pattes I2C non configuré.\\
Il faut l'importer via le module de Recherche BOARD de PlatformIO, puis chercher le nom de la carte.\\

On se renseigne sur le RAK11300 utilisé en WisBlock via son Datasheet.\\
Pour le connector A, nous sommes sur du I2C.\\

Nous avons séléctionné la bonne carte, mais nous cherchons le fichier de définition des bords par WireBoard.\\ --> Il faut faire un patch ?\\

\href{https://github.com/RAKWireless/WisBlock/tree/master/PlatformIO} {Github RAKWireless}

\section{Conclusion}

\section{Bonus : Journal de Bord}

\subsection{15 Septembre 2026}

\textit{Nous avons seulement installé 8 vis, 4 pour chaque puce et avons branché un fil en I2C : le capteur de température.}\\

\subsection{29/09/2026}

\textit
{
  Le Bootloading via double TAP rapide sur Reset, peut être dangereux : il ne permet pas de drag and drop.\\
  On utilise l'extension GitHub PlatformIO avec le module Arduino (Pas commun selon le professeur) donc il faut se préparer à des difficultés.\\
  Après quelques soucis de connexion avec la carte, la communication s'est finalement mise à marcher via l'ajout de Flags !\\
}

\end{document}
